{\scriptsize\setlength{\tabcolsep}{3pt}
\begin{longtable}{l p{7.5cm} p{3cm} p{5cm}}
\caption{Measurement notes and sources for the observed and agency-reported observations (numbers as in Table~\ref{tab:perf}). The design-calculated rows come from CSE model-project pages, \url{https://www.cseindia.org/rainwater-harvesting-case-studies}, with volumes stated as a fixed share of area $\times$ rainfall, and from the CGWB design note for Shram Shakti Bhawan.}\label{tab:perfnotes}\\
\toprule
\textbf{No.} & \textbf{Measurement note} & \textbf{Source} & \textbf{URL}\\ \midrule\endfirsthead
\toprule \textbf{No.} & \textbf{Measurement note} & \textbf{Source} & \textbf{URL}\\ \midrule\endhead
\bottomrule\endfoot
1 & Flow meter in line. CGWB: 215.7 m3 received on roof, 19 m3 recharged (as per estimation). Rainfall = 215.7/360. & CGWB 2011, p.73 & \url{https://cgwb.gov.in/cgwbpnm/download/379}\\
2 & Flow meter in line (same CGWB wording as Ramtek). Volume stated as 280.17 m3 (90\%) of 311.30 m3 received; the round 90\% may be an assumption, see Results. Rainfall = 311.3/500. & CGWB 2011, p.73 & \url{https://cgwb.gov.in/cgwbpnm/download/379}\\
3 & CGWB: about 830 m3 recharged from 1660 m2 roof; period not stated. Rainfall = IMD Safdarjung normal 1991-2020. & CGWB 2011, p.44 & \url{https://cgwb.gov.in/cgwbpnm/download/379}\\
4 & Recharge measured with a device at the bore well (CSE); Chachadi 2013 gives 934 kL to 11 Nov 2010. 2010 monsoon rainfall 3.7 m reported. & CSE 2014 Urban RWH report; Chachadi 2013 & \url{https://cdn.cseindia.org/userfiles/Urban%20Rainwater%20harvesting%20report.pdf}\\
5 & System built July 2008, so 2008 is a partial season. Rainfall = IMD Panaji normal 1981-2010. & CSE 2014 Urban RWH report & \url{https://cdn.cseindia.org/userfiles/Urban%20Rainwater%20harvesting%20report.pdf}\\
6 & Rainfall = IMD Panaji normal 1981-2010 (2009 site rainfall not reported). & CSE 2014 Urban RWH report & \url{https://cdn.cseindia.org/userfiles/Urban%20Rainwater%20harvesting%20report.pdf}\\
7 & 'Around 1100 KL of rainwater harvested annually' for DM water production; method not stated. Rainfall = IMD Bengaluru normal 1991-2020. & NWM/BEL case sheet & \url{https://www.nwm.gov.in/sites/default/files/BEL_RWH-Ctr.pdf}\\
8 & Lok Sabha reply: 1985 m3 (93\% of total generated runoff) recharged annually. Roof area from CSE 2010. No IMD normal located for Chandigarh; CSE city figure used. & Lok Sabha SQ 384 (23.08.2004); CSE 2010 & \url{https://eparlib.sansad.in/bitstream/123456789/520443/1/2310.pdf}\\
9 & Lok Sabha reply: 3812 m3 recharged annually. Roof area 3550 m2 from CSE 2010. & Lok Sabha SQ 384 (23.08.2004); CSE 2010 & \url{https://eparlib.sansad.in/bitstream/123456789/520443/1/2310.pdf}\\
10 & CGWB Table 34; project 'monitored weekly'; method of volume not stated; 2004 value equals 85\% of gross rainfall exactly. & CGWB 2011, p.116 & \url{https://cgwb.gov.in/cgwbpnm/download/379}\\
11 & CGWB Table 34. & CGWB 2011, p.116 & \url{https://cgwb.gov.in/cgwbpnm/download/379}\\
12 & CGWB Table 34. & CGWB 2011, p.116 & \url{https://cgwb.gov.in/cgwbpnm/download/379}\\
13 & CSE: 'collects about 1.05 million litres per year'; method not stated. & CSE case study page & \url{https://www.cseindia.org/node/1099}\\
14 & CWC: 'Water harvested : 10 lakh litres (approx)'; method not stated. & CWC RWH note & \url{https://cwc.gov.in/sites/default/files/rwh.pdf}\\
15 & 3 kL tank fills four to five times a year (12-15 kL); storage-limited, so a lower bound on roof yield. Midpoint used. & CSE case study page & \url{https://www.cseindia.org/node/1099}\\
16 & About 15,000 L collected over 5 months with a 1 kL tank; storage-limited lower bound. & India Water Portal (Ramamohan 2012) & \url{https://oldiwp.indiawaterportal.org/sites/default/files/iwp2/rooftop_rainwater_harvesting_system___hyderabad___rv_ramamohan___2012.pdf}\\
\end{longtable}}
